\documentclass[10pt,a4paper]{amsart}
\usepackage[margin=19mm]{geometry}
\usepackage{amsmath,amssymb,mathtools}
\usepackage[scr=rsfs,cal=euler]{mathalfa}
\usepackage{xfrac}
\usepackage[unicode,hidelinks,bookmarks=false]{hyperref}
\newcommand{\ie}{i.e.\ }
\newcommand{\cf}{cf.\ }
\newcommand{\resp}{respectively\ }
\newcommand{\Subsection}{\subsection}
\providecommand{\nobreakdash}{\nobreak\hbox{-}}
\setlength{\emergencystretch}{2em}
\multlinegap0pt
\usepackage{tikz}
\newtheorem{lemma}{Lemma}
\title{Low-rank approximation of analytic kernels}
\author{Marcus Webb}
\date{Source: arXiv:2509.14017v4 (CC BY 4.0)}
\begin{document}
\maketitle

\noindent\emph{Source and licence.} Low-rank approximation of analytic kernels. Version arXiv:2509.14017v4 (CC BY 4.0), distributed by arXiv under the Creative Commons Attribution 4.0 International licence, http://creativecommons.org/licenses/by/4.0/, as recorded in the arXiv licence metadata for this exact version and shown on the abstract page https://arxiv.org/abs/2509.14017v4. Full licence notice: \texttt{licenses/arxiv-2509.14017-CC-BY-4.0.txt}.\par\smallskip

\noindent\textit{Editorial scope.} Counted unit: the article's lemma lem:boundZ (source0.tex lines 586--599). The article's complete original proof of it is delivered with it (source0.tex lines 704--706, one \texttt{proof} environment of 3 lines, which the article places away from the statement). No mathematics is written, repaired or re-derived: the statement and the proof are the article's own lines, in the article's own notation, and no retained line is substituted or re-flowed.  Retained uncounted as the source-local context the proof needs: lines 685--687.  Nothing was omitted from any retained span, so the package carries no omission record.  No retained line contains a citation, so the wrapper carries no bibliography and no bibliography entry.  No retained line contains a figure environment, an image-inclusion command or drawing code, so no image is delivered and the package declares no asset dependency.  Editorial formatting only, in addition: an AMS \texttt{amsart} preamble that restates the article's own declarations and macros used by the retained lines, namely the environments \texttt{lemma} on separate counters, together with the standard packages those lines need.  The retained environments print in this document as Lemma 1 in order of appearance, whereas the article prints the same items with the numbering of its own full document.  The complete native source of the article as distributed in the arXiv e-print for this exact version is bundled unchanged as \texttt{sources/185209/source0.tex}.
\begin{equation}\label{eqn:KandK'}
    \frac{1}{2\pi \mathrm{i}}\int_\Gamma K(x,\xi) \, h(\xi) \, \mathrm{d}\xi = \int_F K'(x,z) \, h(z) \, \mathrm{d}\nu(z),
\end{equation}

\begin{lemma}\label{lem:boundZ}
Let $Z \in C(F \times E)$ define an operator $\mathcal{Z} : L_\mu^{p^*}(E) \to L_\nu^q(F)$ for $p,q \in [1,\infty]$, where $p^*$ is the H\"older dual exponent, satisfying $\tfrac{1}{p} + \tfrac{1}{p^*} = 1$. Then
\begin{equation*}
    \|\mathcal{Z}\|_{L_\mu^{p^*}(E) \to L_\nu^q(F)} \leq \|Z\|_{L_\mu^p\left(E; \, L_\nu^q(F)\right)}.
\end{equation*}
\begin{proof}
    We can work through this directly from the definition. For $p \in [1,\infty]$ and $q \in [1,\infty)$,
        \begin{eqnarray*}
            \|\mathcal{Z}\|_{L_\mu^{p^*}(E) \to L_\nu^q(F)} &=& \sup_{\|g\|_{p^*}=1} \left(\int_F \left| \int_E Z(z,y)\, g(y) \, \mathrm{d}\mu(y) \right|^q \, \mathrm{d} \nu(z)\right)^{1/q} \\
            &\leq& \sup_{\|g\|_{p^*}=1} \int_E \left(\int_F \left|Z(z,y) \right|^q \, \mathrm{d}\nu(z)  \right)^{1/q} \, |g(y)| \, \mathrm{d}\mu(y) \qquad \text{ (Minkowski's inequality)}.
        \end{eqnarray*}
        The result now follows by H\"older's inequality. For $q = \infty$, the proof is similar.
    \end{proof}
\end{lemma}

    \begin{proof}
This follows from Equation \eqref{eqn:KandK'} and Lemma \ref{lem:boundZ}.
    \end{proof}

\end{document}
